Para el cumplimiento de los dos primeros objetivos específicos del trabajo, se llevó a cabo una revisión de literatura orientada, en primer lugar, a identificar las teorías y enfoques clave relacionadoscon el aprendizaje autónomo y continuo, así como las estrategias efectivas que fortalecen la capacidad de aprender a aprender. En segundo lugar, esta revisión tuvo como propósito investigar las metodologías y técnicas empleadas para fomentar la escucha activa durante las interacciones comunicativas, tomando como referencia la literatura académica y las mejores prácticas en el campo de la comunicación interpersonal.
\\
\\
El análisis bibliográfico consideró investigaciones recientes sobre metacognición, aprendizaje autorregulado, fijación de metas, estrategias de aprendizaje y procesos de autorreflexión, aspectos que permiten comprender de manera integral la competencia de aprender a aprender. Asimismo, se revisaron estudios centrados en la comprensión auditiva, la escucha activa, la atención, la interpretación de mensajes y las dinámicas de comunicación efectiva, los cuales aportan lineamientos para el fortalecimiento de esta habilidad mediante actividades estructuradas y deliberadas.
\\
\\
Adicionalmente, esta revisión permitió establecer la relación directa entre los aportes de literatura y las funcionalidades implementadas en la plataforma. Elementos como la estructura de metas--objetivos--tareas, las autoevaluaciones, el uso de estrategias de seguimiento del tiempo y la incorporación de desafíos auditivos con retroalimentación inmediata se sustentan en marcos teóricos ampliamente documentados.

\section{Metodología de la revisión}

Para esta revisión se empleó una metodología estructurada que incluyó:

\begin{itemize}
    \item \textbf{Búsqueda bibliográfica:} Se consultaron bases de datos académicas reconocidas, entre ellas IEEE Xplore, ACM Digital Library, ScienceDirect, Scopus y Google Scholar. Adicionalmente, se revisaron libros, capítulos y artículos de acceso abierto relacionados con metacognición, aprendizaje autorregulado, estrategias de aprendizaje y comunicación interpersonal.

    \item \textbf{Criterios de inclusión:} Se consideraron artículos publicados entre 2010 y 2024 que abordaran alguno de los siguientes temas: teorías y enfoques sobre aprendizaje autónomo y continuo, modelos y estrategias de aprendizaje autorregulado, fijación de metas, uso de ePortfolios, estrategias de gestión del tiempo, comprensión auditiva, escucha activa, interacción comunicativa y prácticas pedagógicas para el fortalecimiento de habilidades de comunicación. También se incluyeron estudios que presentaran resultados empíricos o propuestas metodológicas aplicables al contexto educativo universitario.

    \item \textbf{Palabras clave:} Para la búsqueda se emplearon combinaciones de términos en español e inglés, tales como: “aprendizaje autorregulado”, “metacognición”, “learn to learn”, “self-regulated learning”, “goal setting strategies”, “ePortfolio for learning”, “active listening”, “listening comprehension”, “communication skills”, “interpersonal communication” y “listening strategies”.

    \item \textbf{Revisión de artículos seleccionados:} Para el componente de aprendizaje autónomo, se analizaron los artículos recopilados en la tabla de estudios sobre \textit{Learn to Learn} (véase \hyperref[anexo:learn_to_learn]{Anexo A}), los cuales abordan temas como autorregulación, fijación de metas, reflexión, motivación, gestión del tiempo y uso de tecnologías para apoyar procesos metacognitivos. En cuanto al componente de escucha activa, se revisaron artículos que abordan el procesamiento cognitivo y metacognitivo de la comprensión auditiva, incluyendo la identificación de ideas clave, el parafraseo, la verificación de la interpretación, la segmentación del discurso, la memoria de trabajo, y el uso de recursos digitales, actividades guiadas y herramientas gamificadas para fortalecer la precisión interpretativa y la atención en tareas auditivas (véase \hyperref[anexo:escucha_activa]{Anexo B}).

    \item \textbf{Análisis comparativo:} Finalmente, se identificó cómo los hallazgos de la literatura se aplican en las funcionalidades implementadas en los módulos de la plataforma, con el propósito de asegurar que cada componente desarrollado se sustentara en enfoques y prácticas respaldadas por la investigación académica.
\end{itemize}

\section{Síntesis de estudios sobre \textit{Learn to Learn}}

La revisión de literatura realizada permitió identificar un conjunto de enfoques, teorías y estrategias que fundamentan el desarrollo del módulo \textit{Learn to Learn}. Los estudios seleccionados abordan componentes esenciales del aprendizaje autónomo y autorregulado, tales como la metacognición, la fijación de metas, la planificación, la gestión del tiempo, la motivación intrínseca y la reflexión continua. Estos elementos han sido ampliamente documentados en la literatura como factores determinantes para fortalecer la capacidad de los estudiantes de gestionar su propio proceso de aprendizaje.
\\
\\
A partir del análisis de los artículos incluidos (véase \hyperref[anexo:learn_to_learn]{Anexo A}), se elaboró una tabla sintetizada que reúne los aportes más significativos de cada estudio y su relación directa con las funcionalidades implementadas en el módulo. Aunque algunos trabajos presentan modelos o fundamentos teóricos, otros ofrecen hallazgos empíricos, propuestas metodológicas o marcos prácticos, razón por la cual la síntesis integra distintos tipos de contribuciones relevantes para el desarrollo del módulo.
\\
\\
Esta síntesis busca mostrar de manera clara cómo cada concepto o estrategia identificada en la literatura contribuye al diseño del modelo de metas--objetivos--tareas, a las prácticas de autorreflexión y a las herramientas de seguimiento utilizadas en la plataforma.
\\
\\
La Tabla~\ref{tab:syntesis_learn_to_learn} reúne los principales estudios analizados para fundamentar el módulo Learn to Learn, organizados según el concepto central abordado en cada artículo, la contribución específica del estudio y su aplicación en los componentes del módulo.

\begin{center}
\footnotesize
\setlength{\tabcolsep}{3pt}
\renewcommand{\arraystretch}{1.2}

    \begin{longtable}{| m{4.2cm} | m{3cm} | m{4.5cm} | m{4.5cm} |}
    \caption{Síntesis de estudios relevantes para el módulo \textit{Learn to Learn}}
    \label{tab:syntesis_learn_to_learn} \\
    \hline
    \textbf{Artículo} & \textbf{Concepto clave} & \textbf{Contribución del estudio} & \textbf{Aplicación en el módulo} \\
    \hline
    \endfirsthead

    \hline
    \textbf{Artículo} & \textbf{Concepto clave} & \textbf{Contribución del estudio} & \textbf{Aplicación en el módulo} \\
    \hline
    \endhead
    
    \hline \multicolumn{4}{c}{{\textit{Continúa en la siguiente página}}} \\
    \endfoot
    
    \hline
    \endlastfoot
    
    \textit{Hacia una nueva conceptualización del Aprender a Aprender} (Caena \& Stringher, 2020) \cite{caena_stringher_2020_aprender} &
    Conceptualización de AaA &
    El estudio redefine el Aprender a Aprender como una hiper-competencia que integra procesos metacognitivos, autorregulatorios y socio-afectivos para dirigir el aprendizaje de manera consciente. Propone un proceso ejecutivo que permite planificar, monitorear y evaluar el propio aprendizaje, impulsado por la autonomía como necesidad psicológica clave. Además, evidencia que el AaA implica ajustarse y responder a distintas tareas, entornos y demandas socioculturales, mostrando que no es una habilidad aislada sino un conjunto articulado de procesos que apoyan la autorregulación y la adaptación del aprendiz &
    Sustenta la definición del módulo como una competencia transversal que integra la fijación de metas, la reflexión y procesos de autorregulación del aprendizaje. \\
    \hline
    
    \textit{Learning how to learn and assessment for learning: a theoretical inquiry} (Black et al., 2007) \cite{black_etal_2007_learning} &
    AaA y evaluación formativa &
    Analiza cómo la evaluación formativa (AfL) contribuye al desarrollo del aprender a aprender al promover la reflexión, el feedback accionable y la toma de conciencia del propio progreso. El estudio muestra que la retroalimentación continua y co-construida fortalece la autonomía del aprendiz, fomenta la intencionalidad en la toma de decisiones y crea un entorno social que facilita la autorregulación y la mejora del desempeño. &
    Fundamenta la integración de autoevaluaciones y retroalimentación continua para promover la reflexión, la autorregulación y la toma de decisiones que guían el progreso del estudiante en el módulo. \\
    \hline
    
    \textit{Becoming a self-regulated learner: an overview} (Zimmerman, 2010) \cite{Zimmerman2002}&
    Aprendizaje autorregulado (SRL) &
    Presenta el modelo cíclico de autorregulación del aprendizaje compuesto por las fases de planificación, ejecución y autorreflexión. El artículo describe procesos como la fijación de metas, el monitoreo estratégico, el control motivacional y la autoevaluación, mostrando cómo estos permiten al estudiante gestionar activamente su aprendizaje. &
    Sustenta la organización del módulo en torno a las fases de planificación, ejecución y reflexión, y respalda el uso de metas, objetivos, tareas y monitoreo del progreso como componentes centrales de la autorregulación. \\
    \hline
    
    \textit{Do students learn what they teach when generating teaching materials for others?} (Ribosa \& Durán, 2022) \cite{ribosa_duran_2022_teach}&
    Aprender enseñando &
    Meta-análisis que demuestra que crear materiales didácticos para otros mejora el aprendizaje del contenido enseñado. Explica que este proceso activa mecanismos cognitivos generativos (como organizar, seleccionar y transformar la información) que promueven una comprensión más profunda y una mejor estructuración del conocimiento. También evidencia que los materiales visuales y audiovisuales potencian más el aprendizaje que los textuales. &
    Se relaciona directamente con el módulo, ya que la construcción de los planes de aprendizaje (roadmaps) (mediante la formulación de metas, objetivos y tareas) exige que el estudiante estructure, organice y represente visualmente su proceso de aprendizaje. Este acto de generar una explicación o material visual activa los mismos procesos cognitivos generativos descritos en el aprendizaje por enseñanza, favoreciendo una comprensión más profunda. \\
    \hline
    
    \textit{Personal OKRs: A tool for academic goal setting} (2023) \cite{lertladaluck_2023_personal_okrs}&
    Fijación de metas (OKRs) &
    El estudio implementa el marco OKR como herramienta para establecer metas académicas y demuestra, mediante una intervención de 16 semanas, que el uso de objetivos claros, resultados clave medibles y revisiones periódicas favorece la calidad del establecimiento de metas. Los hallazgos muestran una mejora significativa en el componente de metas basadas en datos (considerar experiencias previas, habilidades e información externa) impulsada por el seguimiento frecuente y la autorreflexión. &
    Alinea la estructura del módulo Learn to Learn (metas, objetivos, tareas y seguimiento) con los principios del marco OKR. Asimismo, fundamenta la incorporación de procesos de revisión continua, reflexión y ajuste del plan de aprendizaje \\
    \hline
    
    \textit{Implementing a Self-Regulated Oriented ePortfolio} (Alexiou \& Paraskeva, 2014) \cite{alexiou_paraskeva_2014_eportfolio}&
    ePortfolio y SRL &
    Presenta un ePortfolio orientado a la autorregulación construido sobre el modelo cíclico de Zimmerman, donde los estudiantes establecen metas, seleccionan estrategias, registran evidencias y reflexionan sobre su progreso. Los resultados muestran una mejora significativa en la calidad del establecimiento de metas tras su uso, indicando que el ePortfolio fortalece el monitoreo del aprendizaje, la identificación de fortalezas y dificultades, y la capacidad de diseñar estrategias de mejora. &
    Se relaciona con el módulo al fundamentar el uso de registros de progreso, reflexiones guiadas y seguimiento continuo del aprendizaje. Esta correspondencia se refleja particularmente en el formulario de reporte de aprendizaje que el estudiante completa al finalizar una meta, donde evalúa sus avances, dificultades y estrategias, alineándose con las prácticas de autorreflexión promovidas en los ePortfolios orientados a la autorregulación. \\
    \hline
    
    \textit{What Makes Adults Choose to Learn} (2018) \cite{broek_etal_2023_adults_choose}&
    Agencia y motivación &
    Analiza los factores que estimulan o inhiben la decisión de los adultos de aprender, destacando el papel de la motivación intrínseca, la relevancia personal, las emociones, la autoeficacia y el apoyo social. También identifica barreras como la falta de tiempo, experiencias negativas previas y la ausencia de claridad en los objetivos. El estudio subraya que la agencia del aprendiz (la sensación de control, autonomía y alineación con sus metas reales) es determinante para iniciar y sostener procesos de aprendizaje. &
    Justifica la personalización del módulo, ya que el estudio demuestra que la decisión de aprender depende de factores individuales como la motivación, la autonomía, la relevancia personal y las barreras percibidas. La formulación de metas propias, el reconocimiento de dificultades durante la reflexión final de cada meta y la adaptación del plan de aprendizaje (roadmap) responden directamente a las necesidades reales del adulto aprendiz. \\
    \hline
    
    \textit{Supporting students’ goal setting process using chatbot} (Du et al., 2021) \cite{du_etal_2021_chatbot_goals} &
    Apoyo digital a metas &
    El estudio implementa un chatbot basado en el marco SMART para guiar a los estudiantes en la fijación de metas al inicio de un curso en línea. A través de preguntas estructuradas y recomendaciones personalizadas, el chatbot ayudó a los estudiantes a clarificar y visualizar sus objetivos, formular metas más realistas y tomar conciencia del valor de establecer intenciones antes de comenzar a aprender. Los resultados evidencian que esta intervención digital fortalece la planificación inicial dentro del proceso de autorregulación del aprendizaje. &
    El estudio respalda el uso de apoyo digital para guiar la formulación de metas y pasos concretos, lo cual se refleja en el módulo mediante Aithena, el agente IA encargado de generar un plan de aprendizaje inicial a partir de los intereses, conocimientos previos y necesidades del estudiante. \\
    \hline
    
    \textit{Video Game Design for Learning to Learn} (Pasqualotto et al., 2023) \cite{pasqualotto_2023_videogame_learn2learn} &
    Diseño gamificado &
    El estudio analiza principios de diseño de videojuegos que favorecen el aprendizaje, destacando la importancia de la retroalimentación continua, la progresión por niveles, el ajuste dinámico de la dificultad y la estructuración guiada de la experiencia. Estos elementos mantienen la motivación, orientan la atención hacia los objetivos relevantes y permiten un avance gradual y significativo en las habilidades a desarrollar. &
    Respaldan la estructura progresiva del módulo, en la que el aprendizaje se organiza en pasos guiados (metas → objetivos → tareas), permitiendo que el estudiante avance de manera gradual y consciente, con apoyo estructurado durante la planificación y claridad sobre el progreso en cada etapa. \\
    \hline

    \end{longtable}
\end{center}

\section{Síntesis de estudios sobre Active Listening}
La revisión de literatura realizada para el módulo de \textit{Active Listening} permitió identificar un conjunto de enfoques, prácticas pedagógicas y procesos cognitivos que fundamentan el desarrollo de actividades orientadas al fortalecimiento de la comprensión auditiva. Los estudios seleccionados abordan componentes esenciales de la escucha activa, tales como la atención sostenida, la identificación de ideas clave, el parafraseo, la verificación de la comprensión, la segmentación del discurso, la memoria de trabajo y el monitoreo metacognitivo. Estos elementos han sido ampliamente documentados en la literatura como factores que influyen de manera directa en la precisión interpretativa y en la capacidad de reconstruir significados a partir de información auditiva limitada.
\\
\\
A partir del análisis de los artículos incluidos (véase \hyperref[anexo:active_listening]{Anexo B}), se elaboró una tabla que sintetiza los aportes más relevantes de cada estudio y su relación con las funcionalidades implementadas en el módulo. Mientras algunos trabajos proponen estrategias instruccionales (como actividades guiadas de parafraseo, síntesis o práctica deliberada) otros ofrecen hallazgos empíricos sobre el impacto del procesamiento cognitivo y metacognitivo en la escucha, o bien exploran el papel de tecnologías interactivas, recursos multimedia y herramientas gamificadas para mejorar la motivación, la atención y la claridad interpretativa durante tareas auditivas.
\\
\\
Esta síntesis busca evidenciar cómo cada concepto, estrategia o enfoque identificado en la literatura contribuye al diseño de las rondas auditivas del módulo, a los mecanismos de retroalimentación inmediata y a las dinámicas interactivas que favorecen la precisión interpretativa y el monitoreo de la comprensión durante las tareas auditivas.
\\
\\
La Tabla~\ref{tab:syntesis_active_listening} presenta los estudios analizados, organizados según su concepto central, la contribución específica del artículo y su aplicación en los componentes del módulo de escucha activa.

\begin{center}
\footnotesize
\setlength{\tabcolsep}{3pt}
\renewcommand{\arraystretch}{1.2}

\begin{longtable}{| m{4.2cm} | m{3cm} | m{4.5cm} | m{4.5cm} |}
\caption{Síntesis de estudios relevantes para el módulo \textit{Active Listening}}
\label{tab:syntesis_active_listening} \\
\hline
\textbf{Artículo} & \textbf{Concepto clave} & \textbf{Contribución del estudio} & \textbf{Aplicación en el módulo} \\
\hline
\endfirsthead

\hline
\textbf{Artículo} & \textbf{Concepto clave} & \textbf{Contribución del estudio} & \textbf{Aplicación en el módulo} \\
\hline
\endhead

\hline \multicolumn{4}{c}{{\textit{Continúa en la siguiente página}}} \\
\endfoot

\hline
\endlastfoot

\textit{Active Listening for Engineering Students} (Novick, Realyvasquez and Ramirez, 2021) \cite{novick_2021_active_listening}  &
Entrenamiento estructurado en escucha activa &
El estudio identifica la escucha activa como una de las habilidades más demandadas por empleadores de ingenieros y desarrolla un conjunto de módulos formativos distribuidos a lo largo de un programa universitario. El artículo define la escucha activa en cuatro fases (pre-listening, during, post-listening inmediata y a largo plazo) e integra comportamientos concretos como evitar respuestas juiciosas, ofrecer retroalimentación verbal y no verbal adecuada, mantener la atención y formular preguntas para fomentar la expresión del interlocutor. Además, describe actividades pedagógicas efectivas (como roles de escucha, ejercicios de práctica deliberada, prompts guiados y reflexión posterior) que permiten a los estudiantes comprender, practicar y evaluar habilidades de escucha activa. &
Sustenta la inclusión de actividades auditivas estructuradas, dado que el estudio evidencia que la práctica guiada, el modelamiento y la retroalimentación inmediata fortalecen comportamientos esenciales de escucha activa. \\ \hline

\textit{Active Listening in Written Online Communication} (Bauer and Figl, 2008)  \cite{bauer_figl_2008_active_listening_online} &
Escucha activa en entornos sin señales no verbales &
El estudio analiza la aplicabilidad de las técnicas de escucha activa en comunicación escrita mediante mensajería instantánea, un contexto que carece de claves no verbales y presenta mayor riesgo de distracciones. Los resultados muestran que todas las técnicas clásicas de escucha activa (parafrasear, clarificar, preguntar, resumir y animar al interlocutor) aparecieron de manera consistente en las interacciones analizadas, demostrando que pueden usarse eficazmente incluso en entornos textuales. Sin embargo, el estudio identifica desafíos importantes: mensajes superpuestos que inducen malentendidos, distracciones frecuentes, dificultad para mantener la atención y limitaciones para expresar o interpretar emociones sin apoyo vocal o gestual. También se reconoce como ventaja la permanencia del texto, que permite revisar mensajes y apoyar la comprensión. &
El estudio respalda el diseño de actividades donde el estudiante debe comprender un mensaje exclusivamente a partir del contenido verbal, sin apoyo visual ni claves no verbales, tal como ocurre en las rondas de audio del módulo. Además, el artículo evidencia que la falta de atención sostenida y los malentendidos son problemas frecuentes en contextos sin claves no verbales, lo que respalda el uso de rondas secuenciales, repeticiones controladas del audio y retroalimentación inmediata en el módulo, ya que estos elementos ayudan a reforzar la comprensión y reducir errores interpretativos. \\ \hline

\textit{“Can You Repeat That?” Teaching Active Listening in Management Education} (Spataro and Bloch, 2018) \cite{spataro_bloch_2018_active_listening_management} &
Pedagogías de parafraseo, síntesis y práctica guiada &
El estudio propone un método pedagógico para enseñar escucha activa basado en autoevaluación, materiales instruccionales, práctica deliberada y reflexión estructurada. Demuestra que actividades guiadas que incluyen parafraseo, formulación de preguntas, síntesis y verificación de comprensión permiten a los estudiantes interiorizar las habilidades de escucha activa. Los resultados cuantitativos y cualitativos evidencian mejoras en la capacidad de describir estrategias de escucha, reconocer cuándo aplicarlas y transferirlas a nuevos contextos. &
Sustenta el uso de actividades que requieren parafrasear, sintetizar, completar información, identificar ideas centrales y formular preguntas de aclaración, tal como ocurre en los modos de juego Paraphrase, Summarize, Focus, Cloze y Clarify del módulo. El énfasis del estudio en la práctica deliberada y en ejercicios guiados respalda la estructura por rondas, la repetición controlada del audio y la retroalimentación inmediata, que permiten al estudiante internalizar gradualmente las habilidades de escucha activa. \\ \hline

\textit{The Relative Effectiveness of Active Listening in Initial Interactions} (Weger, Castle and Emmett, 2014) \cite{weger_castle_emmett_2014_relative_effectiveness} &
Efectos del parafraseo y retroalimentación empática &
El estudio compara el parafraseo con otros estilos de respuesta en interacciones iniciales y demuestra que las respuestas basadas en escucha activa incrementan significativamente la percepción de empatía y apoyo. Sin embargo, también evidencia que el parafraseo no garantiza que el hablante se sienta completamente comprendido, mostrando que la interpretación del mensaje puede diferir de la intención original. Estos hallazgos resaltan la complejidad de lograr una comprensión precisa incluso cuando se aplican técnicas de escucha activa. &
Refuerza el valor de actividades donde el estudiante compara su interpretación del mensaje con el contenido real, como ocurre en los modos de juego Paraphrase, Summarize, Focus y Clarify. El hallazgo de que el parafraseo no siempre conduce a una comprensión completa justifica el uso de retroalimentación inmediata en el módulo, permitiendo que el estudiante verifique si lo que entendió coincide con el mensaje original y ajuste progresivamente su precisión interpretativa a lo largo de las rondas. \\ \hline

\textit{Listening Strategy Applications by Learners under the Context of Multimodality} (Kao, Liao and Lan, 2017) \cite{kao_liao_lan_2017_listening_strategies} &
Estrategias de escucha y uso de estrategias metacognitivas bajo modalidades multimodales &
El estudio aporta evidencia empírica sobre cómo los estudiantes gestionan su proceso de comprensión auditiva mediante estrategias metacognitivas (tales como atención selectiva, atención dirigida, identificación de problemas y monitoreo) incluso cuando cambian las condiciones de apoyo (entorno multimodal 3D o ayudas impresas). La contribución central radica en demostrar que la regulación metacognitiva es un componente dominante y estable del proceso de escucha, más influyente que las variaciones en el tipo de apoyo visual disponible. &
Sustenta actividades del módulo que requieren regulación atencional, monitoreo de la comprensión y selección de información relevante. Modos como Focus, Summarize, Paraphrase y Cloze demandan atención dirigida, identificación de ideas clave y verificación de la interpretación, procesos alineados con las estrategias metacognitivas identificadas como predominantes en el estudio. \\ \hline

\textit{Promoting Listening Skills with an Educational Tool: VoScreen} (Merzifonluoğlu, 2023) \cite{merzifonluoglu_2023_voscreen} &
Gamificación aplicada a la comprensión auditiva &
El estudio demuestra que VoScreen, una herramienta basada en microvideos con retroalimentación inmediata, produce mejoras estadísticamente significativas en la comprensión auditiva de los estudiantes. El análisis cualitativo destaca beneficios adicionales como mayor motivación, exposición a acentos variados y la percepción de que el formato breve y repetitivo facilita una práctica sostenida. &
Sustenta el enfoque gamificado del módulo al validar el uso de desafíos auditivos cortos, retroalimentación inmediata y tareas breves que aumentan la motivación, la atención y la continuidad en la práctica de escucha. \\ \hline

\textit{Enhancing 21st Century Students’ Listening Skills via Augmented Reality and Mobile Applications} (Ismayatim et al., 2019) \cite{ismayatim_etal_2019_ar_listening} &
Tecnologías móviles y realidad aumentada aplicadas a prácticas de escucha &
El estudio analiza las percepciones de estudiantes universitarios sobre MyEVO, un módulo de escucha que integra realidad aumentada, dispositivos móviles y videos auténticos. Los resultados muestran una valoración ampliamente positiva: los estudiantes consideran que la combinación de audio y video facilita la comprensión del contenido, incrementa el interés y permite un aprendizaje más flexible e independiente. Asimismo, destacan que las prácticas basadas en AR son interactivas, modernas y fáciles de usar, y que la posibilidad de controlar el ritmo (pausar, repetir, revisar) favorece la experiencia de aprendizaje. &
Apoya el diseño del módulo al mostrar que entornos digitales interactivos que combinan estímulos auditivos con una interfaz multimedia incrementan el interés, la percepción de comprensión y la autonomía del estudiante, según las valoraciones reportadas. En nuestro caso, estos hallazgos respaldan el uso de audios breves integrados en una interfaz gamificada que permite controlar el ritmo de la actividad, repetir el contenido y participar en prácticas de escucha más dinámicas y accesibles. \\ \hline

\textit{La escucha activa y la memorización en la formación de intérpretes} (Bertozzi Díaz, 2021) \cite{bertozzi_2021_escucha_activa_interpretes} &
Memoria de trabajo, análisis del discurso y reconstrucción del mensaje &
El estudio presenta una experiencia didáctica centrada en la escucha activa y las técnicas de memorización para intérpretes en formación. Expone que la escucha activa implica operaciones cognitivas profundas como analizar la estructura del discurso, identificar ideas principales, distinguir información esencial de la redundante y mantener la atención sostenida. Asimismo, muestra que la memorización eficaz se basa en la segmentación del mensaje (chunking), la visualización mental, la identificación de palabras clave y la reconstrucción del contenido a partir de unidades significativas. Los resultados preliminares indican mejoras percibidas en reformulación, priorización y retención de discursos estructurados y no estructurados. &
Sustenta actividades del módulo que exigen retener información auditiva breve, reconstruir el mensaje, identificar patrones y sintetizar ideas. Apoya especialmente modos como Summarize, Paraphrase y Cloze, que requieren análisis del contenido, comprensión profunda y uso activo de la memoria de trabajo para reorganizar e interpretar fragmentos auditivos limitados. \\ \hline

\end{longtable}
\end{center}
